%%%%%%%%%%%%%%%%%%%%%%%%%%%%%%%%%%%%%%%%%%%%%%%%%%%%%%%%%%%%%%%%%%%%%%
% Lớp 9 | Chương V. Đường tròn | Vị trí tương đối của hai đường tròn
% Dạng 9H5?17-1: Xác định vị trí tương đối của hai đường tròn
%%%%%%%%%%%%%%%%%%%%%%%%%%%%%%%%%%%%%%%%%%%%%%%%%%%%%%%%%%%%%%%%%%%%%%


% nguồn: 2026-11-21_lop9_chuong5_bai14_17 (nhập 21/11/2026)
\begin{ex}%[9H5N17-1]%[ma:d070cb53]
% nguon: SGK Toán 9 KNTT tập 1 – Bài 17, Vị trí tương đối của hai đường tròn
	Cho hai đường tròn $(O; 7\text{ cm})$ và $(O'; 3\text{ cm})$ có khoảng cách hai tâm $OO' = 10\text{ cm}$. Vị trí tương đối của hai đường tròn đó là
	\choice
	{Ở ngoài nhau}
	{\True Tiếp xúc ngoài}
	{Cắt nhau}
	{Tiếp xúc trong}
	\loigiai{
		Ta có $R + r = 7 + 3 = 10\text{ cm} = OO'$, nên hai đường tròn $(O)$ và $(O')$ tiếp xúc ngoài với nhau.
	}
\end{ex}

% nguồn: 2026-11-21_lop9_chuong5_bai14_17 (nhập 21/11/2026)
\begin{ex}%[9H5H17-1]%[ma:7a98286a]
% nguon: SGK Toán 9 KNTT tập 1 – Bài 17, Vị trí tương đối của hai đường tròn
	Cho hai đường tròn $(O; 8\text{ cm})$ và $(O'; 5\text{ cm})$ có khoảng cách giữa hai tâm là $OO' = d$. Xét tính đúng, sai của các khẳng định sau:
	\choiceTF
	{\True Nếu $d = 15\text{ cm}$ thì hai đường tròn ở ngoài nhau}
	{\True Nếu $d = 3\text{ cm}$ thì hai đường tròn tiếp xúc trong}
	{Nếu $d = 2\text{ cm}$ thì hai đường tròn cắt nhau}
	{\True Hai đường tròn cắt nhau khi và chỉ khi $3 < d < 13$}
	\loigiai{
		\begin{itemchoice}
			\itemch Đúng. Ta có $R + r = 8 + 5 = 13\text{ cm}$. Vì $d = 15 > 13$ nên hai đường tròn ở ngoài nhau.
			\itemch Đúng. Ta có $R - r = 8 - 5 = 3\text{ cm}$. Vì $d = 3 = R - r$ nên hai đường tròn tiếp xúc trong.
			\itemch Sai. Vì $d = 2 < R - r = 3$ nên hai đường tròn đựng nhau (không giao nhau).
			\itemch Đúng. Hai đường tròn cắt nhau khi và chỉ khi $R - r < d < R + r$, tức là $3 < d < 13$.
		\end{itemchoice}
	}
\end{ex}

% nguồn: 2026-11-21_lop9_chuong5_bai14_17 (nhập 21/11/2026)
\begin{bt}%[9H5V17-1]%[ma:42104565]
% nguon: SGK Toán 9 KNTT tập 1 – Bài 17, Vị trí tương đối của hai đường tròn
	Cho hai đường tròn $(O; 10\text{ cm})$ và $(O'; 17\text{ cm})$ cắt nhau tại hai điểm $A, B$ sao cho dây chung $AB = 16\text{ cm}$ và hai tâm $O, O'$ nằm về hai phía của dây $AB$. Tính khoảng cách $OO'$ giữa hai tâm (đơn vị cm).
	\shortans{21}
	\loigiai{
		Vì $OO'$ là đường trung trực của dây chung $AB$ nên gọi $H = OO' \cap AB$, ta có $OO' \perp AB$ tại $H$ và $AH = \dfrac{AB}{2} = 8\text{ cm}$.\\
		Trong tam giác $OHA$ vuông tại $H$: $OH = \sqrt{OA^2 - AH^2} = \sqrt{10^2 - 8^2} = 6\text{ cm}$.\\
		Trong tam giác $O'HA$ vuông tại $H$: $O'H = \sqrt{O'A^2 - AH^2} = \sqrt{17^2 - 8^2} = 15\text{ cm}$.\\
		Vì $O$ và $O'$ nằm về hai phía của dây $AB$ nên $OO' = OH + O'H = 6 + 15 = 21\text{ cm}$.
	}
\end{bt}

% nguồn: 2026-11-21_lop9_chuong5_bai16_17_bs (nhập 21/11/2026)
\begin{ex}%[9H5N17-1]%[ma:c5b87ba1]
% nguon: SGK Toán 9 KNTT tập 1 – Bài 17, Vị trí tương đối của hai đường tròn
	Cho hai đường tròn $(O; 9\text{ cm})$ và $(O'; 4\text{ cm})$ có khoảng cách giữa hai tâm là $OO' = 5\text{ cm}$. Vị trí tương đối của hai đường tròn đó là
	\choice
	{Ở ngoài nhau}
	{Tiếp xúc ngoài}
	{Cắt nhau}
	{\True Tiếp xúc trong}
	\loigiai{
		Ta có $R - r = 9 - 4 = 5\text{ cm} = OO'$, do đó hai đường tròn tiếp xúc trong với nhau.
	}
\end{ex}

% nguồn: 2026-11-21_lop9_chuong5_bai16_17_bs (nhập 21/11/2026)
\begin{ex}%[9H5N17-1]%[ma:2e0104da]
% nguon: SGK Toán 9 KNTT tập 1 – Bài 17, Vị trí tương đối của hai đường tròn
	Cho hai đường tròn $(O; 6\text{ cm})$ và $(O'; 4\text{ cm})$ có khoảng cách giữa hai tâm là $OO' = 12\text{ cm}$. Vị trí tương đối của hai đường tròn đó là
	\choice
	{\True Ở ngoài nhau}
	{Tiếp xúc ngoài}
	{Cắt nhau}
	{Đựng nhau}
	\loigiai{
		Ta có $OO' = 12\text{ cm} > R + r = 6 + 4 = 10\text{ cm}$, nên hai đường tròn ở ngoài nhau.
	}
\end{ex}

% nguồn: 2026-11-21_lop9_chuong5_bai16_17_bs (nhập 21/11/2026)
\begin{ex}%[9H5H17-1]%[ma:6b06eb8d]
% nguon: SGK Toán 9 KNTT tập 1 – Bài 17, Vị trí tương đối của hai đường tròn
	Cho hai đường tròn $(O; 10\text{ cm})$ và $(O'; 6\text{ cm})$ có khoảng cách giữa hai tâm là $OO' = d$. Xét tính đúng, sai của các khẳng định sau:
	\choiceTF
	{\True Hai đường tròn tiếp xúc ngoài khi $d = 16\text{ cm}$}
	{\True Hai đường tròn tiếp xúc trong khi $d = 4\text{ cm}$}
	{Nếu $d = 10\text{ cm}$ thì hai đường tròn ở ngoài nhau}
	{\True Có đúng $11$ giá trị nguyên của $d$ (tính theo cm) để hai đường tròn cắt nhau}
	\loigiai{
		\begin{itemchoice}
			\itemch Đúng. Hai đường tròn tiếp xúc ngoài khi $d = R + r = 10 + 6 = 16\text{ cm}$.
			\itemch Đúng. Hai đường tròn tiếp xúc trong khi $d = R - r = 10 - 6 = 4\text{ cm}$.
			\itemch Sai. Vì $4 < 10 < 16$ nên khi $d = 10\text{ cm}$ hai đường tròn cắt nhau.
			\itemch Đúng. Hai đường tròn cắt nhau khi $4 < d < 16$; các giá trị nguyên của $d$ là $5, 6, \ldots, 15$, gồm $15 - 5 + 1 = 11$ giá trị.
		\end{itemchoice}
	}
\end{ex}

% nguồn: 2026-11-21_lop9_chuong5_bai16_17_bs (nhập 21/11/2026)
\begin{ex}%[9H5H17-1]%[ma:ecd62316]
% nguon: SGK Toán 9 KNTT tập 1 – Bài 17, Vị trí tương đối của hai đường tròn
	Cho hai đường tròn $(O; 5\text{ cm})$ và $(O'; 5\text{ cm})$ có khoảng cách hai tâm $OO' = 6\text{ cm}$ cắt nhau tại hai điểm $A$ và $B$. Xét tính đúng, sai của các khẳng định sau:
	\choiceTF
	{\True Hai đường tròn cắt nhau tại hai điểm phân biệt}
	{\True Đường nối tâm $OO'$ là đường trung trực của dây chung $AB$}
	{Độ dài dây chung $AB$ bằng $6\text{ cm}$}
	{\True Tứ giác $OAO'B$ là hình thoi có diện tích bằng $24\text{ cm}^2$}
	\loigiai{
		\begin{itemchoice}
			\itemch Đúng. Ta có $5 - 5 = 0 < 6 < 5 + 5 = 10$ nên hai đường tròn cắt nhau tại hai điểm phân biệt.
			\itemch Đúng. Vì $OA = OB = 5\text{ cm}$ và $O'A = O'B = 5\text{ cm}$ nên $OO'$ là đường trung trực của $AB$.
			\itemch Sai. Gọi $H = OO' \cap AB$, ta có $OH = \dfrac{OO'}{2} = 3\text{ cm}$ nên $AH = \sqrt{5^2 - 3^2} = 4\text{ cm}$, suy ra $AB = 8\text{ cm}$.
			\itemch Đúng. Vì $OA = AO' = O'B = BO = 5\text{ cm}$ nên $OAO'B$ là hình thoi, có diện tích $\dfrac{1}{2} \cdot OO' \cdot AB = \dfrac{1}{2} \cdot 6 \cdot 8 = 24\text{ cm}^2$.
		\end{itemchoice}
	}
\end{ex}

% nguồn: 2026-11-21_lop9_chuong5_bai16_17_bs (nhập 21/11/2026)
\begin{bt}%[9H5V17-1]%[ma:10c21510]
% nguon: SGK Toán 9 KNTT tập 1 – Bài 17, Vị trí tương đối của hai đường tròn
	Cho hai đường tròn $(O; 13\text{ cm})$ và $(O'; 15\text{ cm})$ cắt nhau tại $A, B$ sao cho dây chung $AB = 24\text{ cm}$ và hai tâm $O, O'$ nằm về hai phía đối với dây $AB$. Tính khoảng cách $OO'$ giữa hai tâm (đơn vị cm).
	\shortans{14}
	\loigiai{
		Gọi $H = OO' \cap AB$, ta có $OO' \perp AB$ tại $H$ và $AH = \dfrac{AB}{2} = 12\text{ cm}$.\\
		Trong tam giác $OHA$ vuông tại $H$: $OH = \sqrt{13^2 - 12^2} = 5\text{ cm}$.\\
		Trong tam giác $O'HA$ vuông tại $H$: $O'H = \sqrt{15^2 - 12^2} = 9\text{ cm}$.\\
		Vì $O$ và $O'$ nằm về hai phía của dây $AB$ nên $OO' = OH + O'H = 5 + 9 = 14\text{ cm}$.
	}
\end{bt}

% nguồn: 2026-11-21_lop9_chuong5_bai16_17_bs (nhập 21/11/2026)
\begin{bt}%[9H5V17-1]%[ma:7a9c25f1]
% nguon: SGK Toán 9 KNTT tập 1 – Bài 17, Vị trí tương đối của hai đường tròn
	Cho hai đường tròn $(O; 13\text{ cm})$ và $(O'; 15\text{ cm})$ cắt nhau tại $A, B$ sao cho dây chung $AB = 24\text{ cm}$ và hai tâm $O, O'$ nằm về cùng một phía đối với dây $AB$. Tính khoảng cách $OO'$ giữa hai tâm (đơn vị cm).
	\shortans{4}
	\loigiai{
		Gọi $H$ là giao điểm của đường thẳng $OO'$ với dây $AB$, ta có $AH = \dfrac{AB}{2} = 12\text{ cm}$.\\
		Tính được $OH = \sqrt{13^2 - 12^2} = 5\text{ cm}$ và $O'H = \sqrt{15^2 - 12^2} = 9\text{ cm}$.\\
		Vì $O$ và $O'$ nằm cùng một phía đối với $AB$ nên $OO' = |O'H - OH| = 9 - 5 = 4\text{ cm}$.
	}
\end{bt}
